\uuid{cyM7}
\titre{Somme de deux dés}
\auteur{worms}
\organisation{exo7}
\licence{CC BY-NC-SA 4.0}
\datecreate{2017-08-28}
\niveau{L2}
\module{Probabilité et statistique}
\chapitre{Probabilité discrète}
\sousChapitre{Variable aléatoire discrète}
\difficulte{}
\theme{Variables discrètes}
\exo7qcmid{worms.worms.029}
\serie{ordonnee}

\contenu{
\texte{
On considère la variable aléatoire $X$ égale au résultat de la somme de 2 dés à 6 faces équilibrés.
}

\begin{question}
\qclassification{260.03}
\qkeeporder
Quelle est la loi de $X$ ?
\begin{answers}
\bad{$X$ suit la loi uniforme sur $\{1,2,\ldots,12\}$.}
\bad{$X$ suit la loi uniforme sur $\{2,3,\ldots,12\}$.}
\bad{$X$ suit la loi binomiale de paramètres 12 et $1/6$.}
\good{Aucune des réponses ci-dessus.}
\end{answers}
\begin{explanations}
C'est bien sûr la loi sur $\{2,3,...,12\}$ dont les poids sont (exo facile) :

 $\{1/36,2/36,3/36,4/36,5/36,6/36,5/36,4/36,3/36,2/36,1/36\}$.

Les questions qui suivent en découlent facilement.
\end{explanations}
\end{question}

\begin{question}
\qclassification{260.03}
\qkeeporder
Que valent l'espérance et la médiane de $X$ ?
\begin{answers}
\bad{L'espérance de $X$ vaut $6.5$ et la médiane vaut $7$.}
\good{Elles valent toutes les deux 7.}
\bad{Elles valent d'autres valeurs que celles proposées ci-dessus.}
\end{answers}
\end{question}

\begin{question}
\qclassification{260.03}
Que valent $\mathbb{P}(X\geq 4)$ et la variance de $X$ ?
\begin{answers}
\bad{$\mathbb{P}(X\geq 4)=3/4$ et $\mathbb{V}ar(X)=329/6$.}
\bad{$\mathbb{P}(X\geq 4)=1/4$ et $\mathbb{V}ar(X)=35/6$.}
\bad{$\mathbb{P}(X\geq 4)=14/36$ et $\mathbb{V}ar(X)=35/6$.}
\good{$\mathbb{P}(X\geq 4)=11/12$ et $\mathbb{V}ar(X)\simeq 5.8$.}

\end{answers}
%\begin{explanations}
%%% réponse=1-(1+2)/36=11/12   et var=329/6 - 7^2 = 35/6 =5.83333
%\end{explanations}
\end{question}
}
